% GENERATED FILE. Edit canonical lessons, then run npm run generate:books.
% canonical-source-hash: fnv1a64:7ac0f35aaaea40f9
% canonical-lessons: JA-C389-oudouri, JA-C389-roji, JA-C389-uramichi, JA-C389-hodoukyou, JA-C389-nukeru, JA-R389-review-384-389

\chapter{The Main Street and the Alleys}
\label{ch:ja-the-main-street-and-the-alleys}
\hlchaptermodality{389}

\begin{chapteropening}
\textbf{By the end of this chapter:} \emph{I can find my way along streets and alleys.}

Complete the last lesson of chapter 389: I can find my way along streets and alleys.
\end{chapteropening}

\section[\texorpdfstring{\ja{おおどおり} (ōdōri)}{おおどおり (ōdōri)}]{\ja{おおどおり} (ōdōri) — a main street; an avenue}

\label{lesson:JA-C389-oudouri}



\begin{quote}
Before the new one: say the Japanese for free of charge, then the Japanese for requiring a fee.
\end{quote}

\subsection*{You'll want to know: \ja{おおどおり}}
\textbf{\ja{おおどおり}} — \emph{ōdōri} — \textquotedblleft{}a main street; an avenue\textquotedblright{}.

\begin{quote}
\textbf{\ja{おおどおりに} \ja{ひゃっかてんが} \ja{あります。}} — \emph{ōdōri ni hyakkaten ga arimasu} — \textquotedblleft{}There is a department store on the main street.\textquotedblright{}
\end{quote}

\begin{cousinweb}
\textbf{\ja{おおどおり}} (\emph{ōdōri}) is a native Japanese word (wago), inherited rather than borrowed from Chinese, built from \emph{ō}, \textquotedblleft{}big\textquotedblright{} and \emph{dōri}, \textquotedblleft{}street\textquotedblright{} (from \emph{tōru}); its kanji were attached later as kun'yomi (JMdict).
\end{cousinweb}

\begin{grammarlens}[title={What you've built}]
One more word for getting around.
\end{grammarlens}

\subsection*{Your turn}
\begin{itemize}
\raggedright
  \item \emph{Say it:} \emph{ōdōri}
  \item \emph{Say it:} \emph{ōdōri}, once more
  \item \emph{Say it:} \emph{ōdōri ni hyakkaten ga arimasu} — There is a department store on the main street.
  \item \emph{Recall:} say the Japanese for free of charge, then the Japanese for requiring a fee, then say \emph{ōdōri} again
\end{itemize}

\begin{tcolorbox}[breakable,colback=teal!4,colframe=teal!35!black,title={Before you move on}]
What does \ja{おおどおり} mean? (\textquotedblleft{}A main street; an avenue\textquotedblright{}.) Say it once more.
\end{tcolorbox}
\section[\texorpdfstring{\ja{ろじ} (roji)}{ろじ (roji)}]{\ja{ろじ} (roji) — an alley; a lane}

\label{lesson:JA-C389-roji}



\begin{quote}
Before the new one: say the Japanese for requiring a fee, then the Japanese for a main street.
\end{quote}

\subsection*{You'll want to know: \ja{ろじ}}
\textbf{\ja{ろじ}} — \emph{roji} — \textquotedblleft{}an alley; a lane\textquotedblright{}.

\begin{quote}
\textbf{\ja{ろじに} \ja{ちいさい} \ja{みせが} \ja{あります。}} — \emph{roji ni chiisai mise ga arimasu} — \textquotedblleft{}There is a small shop in the alley.\textquotedblright{}
\end{quote}

\begin{cousinweb}
\textbf{\ja{ろじ}} (\emph{roji}) is Sino-Japanese (kango): two Chinese characters read with their on'yomi, readings borrowed from Chinese. They are \emph{ro}, \textquotedblleft{}path\textquotedblright{} (Middle Chinese \emph{luH}) and \emph{ji}, \textquotedblleft{}ground\textquotedblright{} (\emph{dijH}) (KANJIDIC2; Middle Chinese in Baxter's transcription).
\end{cousinweb}

\begin{grammarlens}[title={What you've built}]
One more word for getting around.
\end{grammarlens}

\subsection*{Your turn}
\begin{itemize}
\raggedright
  \item \emph{Say it:} \emph{roji}
  \item \emph{Say it:} \emph{roji}, once more
  \item \emph{Say it:} \emph{roji ni chiisai mise ga arimasu} — There is a small shop in the alley.
  \item \emph{Recall:} say the Japanese for requiring a fee, then the Japanese for a main street, then say \emph{roji} again
\end{itemize}

\begin{tcolorbox}[breakable,colback=teal!4,colframe=teal!35!black,title={Before you move on}]
What does \ja{ろじ} mean? (\textquotedblleft{}An alley; a lane\textquotedblright{}.) Say it once more.
\end{tcolorbox}
\section[\texorpdfstring{\ja{うらみち} (uramichi)}{うらみち (uramichi)}]{\ja{うらみち} (uramichi) — a back street; a byway}

\label{lesson:JA-C389-uramichi}



\begin{quote}
Before the new one: say the Japanese for a main street, then the Japanese for an alley.
\end{quote}

\subsection*{You'll want to know: \ja{うらみち}}
\textbf{\ja{うらみち}} — \emph{uramichi} — \textquotedblleft{}a back street; a byway\textquotedblright{}.

\begin{quote}
\textbf{\ja{うらみちを} \ja{とおって、えきに} \ja{いきました。}} — \emph{uramichi o tōtte, eki ni ikimashita} — \textquotedblleft{}I went to the station by the back streets.\textquotedblright{}
\end{quote}

\begin{cousinweb}
\textbf{\ja{うらみち}} (\emph{uramichi}) is a native Japanese word (wago), inherited rather than borrowed from Chinese, built from \emph{ura}, \textquotedblleft{}back\textquotedblright{} and \emph{michi}, \textquotedblleft{}road\textquotedblright{}; its kanji were attached later as kun'yomi (JMdict).
\end{cousinweb}

\begin{grammarlens}[title={What you've built}]
One more word for getting around.
\end{grammarlens}

\subsection*{Your turn}
\begin{itemize}
\raggedright
  \item \emph{Say it:} \emph{uramichi}
  \item \emph{Say it:} \emph{uramichi}, once more
  \item \emph{Say it:} \emph{uramichi o tōtte, eki ni ikimashita} — I went to the station by the back streets.
  \item \emph{Recall:} say the Japanese for a main street, then the Japanese for an alley, then say \emph{uramichi} again
\end{itemize}

\begin{tcolorbox}[breakable,colback=teal!4,colframe=teal!35!black,title={Before you move on}]
What does \ja{うらみち} mean? (\textquotedblleft{}A back street; a byway\textquotedblright{}.) Say it once more.
\end{tcolorbox}
\section[\texorpdfstring{\ja{ほどうきょう} (hodōkyō)}{ほどうきょう (hodōkyō)}]{\ja{ほどうきょう} (hodōkyō) — a footbridge (over a road)}

\label{lesson:JA-C389-hodoukyou}



\begin{quote}
Before the new one: say the Japanese for an alley, then the Japanese for a back street.
\end{quote}

\subsection*{You'll want to know: \ja{ほどうきょう}}
\textbf{\ja{ほどうきょう}} — \emph{hodōkyō} — \textquotedblleft{}a footbridge (over a road)\textquotedblright{}.

\begin{quote}
\textbf{\ja{ほどうきょうで} \ja{どうろを} \ja{わたりました。}} — \emph{hodōkyō de dōro o watarimashita} — \textquotedblleft{}I crossed the road by the footbridge.\textquotedblright{}
\end{quote}

\begin{cousinweb}
\textbf{\ja{ほどうきょう}} (\emph{hodōkyō}) is Sino-Japanese (kango): three Chinese characters read with their on'yomi, readings borrowed from Chinese. They are \emph{ho}, \textquotedblleft{}walk\textquotedblright{} (Middle Chinese \emph{buH}), \emph{dō}, \textquotedblleft{}street\textquotedblright{} (\emph{dawX}) and \emph{kyō}, \textquotedblleft{}bridge\textquotedblright{} (\emph{gjew}) (KANJIDIC2; Middle Chinese in Baxter's transcription).
\end{cousinweb}

\begin{grammarlens}[title={What you've built}]
One more word for getting around.
\end{grammarlens}

\subsection*{Your turn}
\begin{itemize}
\raggedright
  \item \emph{Say it:} \emph{hodōkyō}
  \item \emph{Say it:} \emph{hodōkyō}, once more
  \item \emph{Say it:} \emph{hodōkyō de dōro o watarimashita} — I crossed the road by the footbridge.
  \item \emph{Recall:} say the Japanese for an alley, then the Japanese for a back street, then say \emph{hodōkyō} again
\end{itemize}

\begin{tcolorbox}[breakable,colback=teal!4,colframe=teal!35!black,title={Before you move on}]
What does \ja{ほどうきょう} mean? (\textquotedblleft{}A footbridge (over a road)\textquotedblright{}.) Say it once more.
\end{tcolorbox}
\section[\texorpdfstring{\ja{ぬける} (nukeru)}{ぬける (nukeru)}]{\ja{ぬける} (nukeru) — to get through (to the other side); to come out}

\label{lesson:JA-C389-nukeru}



\begin{quote}
Before the new one: say the Japanese for a back street, then the Japanese for a footbridge (over a road).
\end{quote}

\subsection*{You'll want to know: \ja{ぬける}}
\textbf{\ja{ぬける}} — \emph{nukeru} — \textquotedblleft{}to get through (to the other side); to come out\textquotedblright{}.

\begin{quote}
\textbf{\ja{ろじを} \ja{ぬけて、えきに} \ja{つきました。}} — \emph{roji o nukete, eki ni tsukimashita} — \textquotedblleft{}I went through the alley and got to the station.\textquotedblright{}
\end{quote}

Polite \textbf{\ja{ぬけます}}, \emph{nukemasu}; an ichidan verb. Past (polite): \textbf{\ja{くるまが} \ja{ろじを} \ja{ぬけました。}}, \emph{kuruma ga roji o nukemashita}.

\begin{cousinweb}
\textbf{\ja{ぬける}} (\emph{nukeru}) is a native Japanese word (wago), inherited rather than borrowed from Chinese; its kanji was attached later as a kun'yomi (JMdict).
\end{cousinweb}

\begin{grammarlens}[title={What you've built}]
One more word for getting around.
\end{grammarlens}

\subsection*{Your turn}
\begin{itemize}
\raggedright
  \item \emph{Say it:} \emph{nukeru}
  \item \emph{Say it:} \emph{nukeru}, once more
  \item \emph{Say it:} \emph{roji o nukete, eki ni tsukimashita} — I went through the alley and got to the station.
  \item \emph{Recall:} say the Japanese for a back street, then the Japanese for a footbridge (over a road), then say \emph{nukeru} again
\end{itemize}

\begin{tcolorbox}[breakable,colback=teal!4,colframe=teal!35!black,title={Before you move on}]
What does \ja{ぬける} mean? (\textquotedblleft{}To get through (to the other side); to come out\textquotedblright{}.) Say it once more.
\end{tcolorbox}
\section[(dialogue)]{Looking back: from To the Airfield to The Main Street and the Alleys}

\label{lesson:JA-R389-review-384-389}



\begin{quote}
Say the Japanese for an airfield and a flight.
\end{quote}

\subsection*{Your turn}
Say these aloud:

\begin{itemize}
\raggedright
  \item an airfield, a flight, takeoff, landing (of a plane), to set off on a journey
  \item overseas, a foreign country, international, a time difference, to aim for
  \item a passport, customs, entry into a country, departure from a country, to take something out
  \item to go out to meet, to happen to see, boarding (a plane or ship), a plane ticket, a business trip
  \item travel expenses, transport costs, a currency, free of charge, requiring a fee
  \item a main street, an alley, a back street, a footbridge (over a road), to get through
\end{itemize}

\begin{tcolorbox}[breakable,colback=teal!4,colframe=teal!35!black,title={Before you move on}]
An airfield? (\textbf{\ja{ひこうじょう}}, \emph{hikōjō}.) A flight? (\textbf{\ja{びん}}, \emph{bin}.) Takeoff? (\textbf{\ja{りりく}}, \emph{ririku}.) Landing (of a plane)? (\textbf{\ja{ちゃくりく}}, \emph{chakuriku}.) To set off on a journey? (\textbf{\ja{たびだつ}}, \emph{tabidatsu}.) Overseas? (\textbf{\ja{かいがい}}, \emph{kaigai}.) A foreign country? (\textbf{\ja{がいこく}}, \emph{gaikoku}.) International? (\textbf{\ja{こくさい}}, \emph{kokusai}.) A time difference? (\textbf{\ja{じさ}}, \emph{jisa}.) To aim for? (\textbf{\ja{めざす}}, \emph{mezasu}.) A passport? (\textbf{\ja{りょけん}}, \emph{ryoken}.) Customs? (\textbf{\ja{ぜいかん}}, \emph{zeikan}.) Entry into a country? (\textbf{\ja{にゅうこく}}, \emph{nyūkoku}.) Departure from a country? (\textbf{\ja{しゅっこく}}, \emph{shukkoku}.) To take something out? (\textbf{\ja{とりだす}}, \emph{toridasu}.) To go out to meet? (\textbf{\ja{でむかえる}}, \emph{demukaeru}.) To happen to see? (\textbf{\ja{みかける}}, \emph{mikakeru}.) Boarding (a plane or ship)? (\textbf{\ja{とうじょう}}, \emph{tōjō}.) A plane ticket? (\textbf{\ja{こうくうけん}}, \emph{kōkūken}.) A business trip? (\textbf{\ja{しゅっちょう}}, \emph{shutchō}.) Travel expenses? (\textbf{\ja{りょひ}}, \emph{ryohi}.) Transport costs? (\textbf{\ja{こうつうひ}}, \emph{kōtsūhi}.) A currency? (\textbf{\ja{つうか}}, \emph{tsūka}.) Free of charge? (\textbf{\ja{むりょう}}, \emph{muryō}.) Requiring a fee? (\textbf{\ja{ゆうりょう}}, \emph{yūryō}.) A main street? (\textbf{\ja{おおどおり}}, \emph{ōdōri}.) An alley? (\textbf{\ja{ろじ}}, \emph{roji}.) A back street? (\textbf{\ja{うらみち}}, \emph{uramichi}.) A footbridge (over a road)? (\textbf{\ja{ほどうきょう}}, \emph{hodōkyō}.) To get through? (\textbf{\ja{ぬける}}, \emph{nukeru}.)
\end{tcolorbox}
